%!TEX root = TQFTnotes.tex
%%%%%%%%%%%%%%%%%%%%%%%%%%%%%%
\chapter{Cobordism hypothesis}\label{c:cob-hyp}
Main reference: \ocite{lurie-cobordism}; short review paper:
\ocite{freed-cobordism}.

\section{Definition of extended TQFT and cobordism hypothesis}\label{s:cobordism-hypothesis}
In this section, we finally give a definition of an extended TQFT in any
dimension $n$.

Recall that in \seref{s:bordn-highercat} we have defined the $(\infty,n)$
category $\Bord_n$ of $n$ bordisms with corners; to be precise, we have defined
two versions: framed bordism category $\Bordfr_n$ and oriented bordism category
$\Bordor_n$. Moreover, each of these categories has a structure of symmetric
monoidal $(\infty,n)$ category, with the monoidal structure given by the
disjoint union.

We can now give the key definition of these lectures.

\begin{definition}\label{d:xtqft-infty}\index{TQFT!extended framed}
  Let $\calC$ be a symmetric monoidal $(\infty, n)$ category. A framed extended
  $n$-dimensional TQFT is a symmetric monoidal functor
  $$
    Z\colon \Bordfr \to \calC.
  $$
  In a similar way one defines the notion of an oriented TQFT.
\end{definition}

We denote the set of all such functors by
$$
  \Fun^{\otimes}(\Bordfr, \calC)
$$
(respectively,   $\Fun^{\otimes}(\Bordfr, \calC)$). By FIXME, this is itself an
$(\infty, 0)$ category, so we can talk, e.g.,  about isomorphism between two
TQFTs.


A fundamental result in this area is the following result.


\begin{theorem}[Cobordism hypothesis]\label{t:cob-hyp}
Let $\calC$ be a symmetric monoidal $(\infty, n)$ category. Then we have an
equivalence of categories
$$
  \Fun^{\otimes}(\Bordfr, \calC)\simeq \cfd
$$
given by $Z\mapsto Z(\bullet)$, where
$\calC^\fd$ the full subcategory of $\calC$ whose objects are fully dualizable
objects  in $\calC$, and   $\cfd$ is the subcategory obtained from
$\calC^\fd$ by discarding all non-invertible morphisms \textup{(}of all orders\textup{)}. By
definition, $\cfd$ is an $(\infty,0)$ category.
\end{theorem}

Note that here it is crucial that we use the framed version of bordism category.

Before discussing the proof, let us note that it immediately implies the
following.
\begin{corollary}
  An $n$-dimensional framed extended  TQFT is  determined uniquely up to
  isomorphism by the object $Z(\bullet)\in \calC$ which must be fully
  dualizable.
\end{corollary}
This is a very far-reaching generalization of \thref{t:1d-tqft} which proved a
similar statement for 1-dimensional TQFTs.

The cobordism hypothesis has a very long history. It was first stated as a
conjecture by Baez and Dolan \ocite{baez-dolan-tqft} in 1995. In 2008, Jacob
Lurie published a famous paper \ocite{lurie-cobordism}, where he  restated this
conjecture in  the language of $(\infty, n)$ categories and outlined the proof.
The crucial ingredient of the proof is result of Igusa on framed Morse functions
\ocite{igusa}. Since then, many people worked on filling the details of the
proof and extending it to other situations, such as \ocite{ayala-francis2017}.


%%%%%%%%%%%%%%%%%%%%%%%%%%%%%%%%%%%%%%%%%%%%%%%%%%%%%%%%%%
\section{$SO(n)$ action and oriented cobordism hypothesis}\label{s:oriented-cobordism-hypothesis}
The cobordism hypothesis gives a description of framed extended TQFTs. Now, let
us discuss other versions of extended TQFT  --- most importantly, oriented
TQFTs.

Since we have an obvious action of the group $\GL(n,\R)$ on $\R^n$, it gives an
action of $\GL(n,\R)$ on the set of all framings of a given manifold and thus on
the set (or rather $\infty$ category)  $\Fun^{\otimes}(\Bordfr, \calC)$ of all
framed TQFTs. By cobordism hypothesis,   $\Fun^{\otimes}(\Bordfr, \calC)\simeq
\cfd$; thus, we get an action of $\GL(n,\R)$ on the $(\infty, 0)$ category
$\cfd$.

Let us now consider oriented extended TQFTs, i.e. monoidal functors $\Bordor\to
\calC$. Since any framed manifold is automatically oriented, we have a functor
$\Bordfr\to\Bordor$; thus, any oriented theory  gives a framed
  theory. Moreover, since  two framings related by the action of the group
$\GL_+(n)$ give the same orientation, we see that an oriented TQFT gives rise to
a $\GL_+(n)$-invariant framed TQFT. Therefore, it is natural to expect some
relation between $\Fun^{\otimes}(\Bordor, \calC)$ and (properly defined)
invariant subcategory $(\cfd)^{\GL_+(n)}$.

To make these ideas precise, we should begin by defining what we mean by action
of a group on a category and by ``invariant subcategory'' under this action.
Unlike group actions on sets, this is far from  obvious. To give the reader some
taste of this, here is the simplest example.

\begin{example}\label{x:group-action-1cat}
  Let $\calC$ be a category (usual 1-category) and let $G$ be a discrete group.
  We define an action of $G$ on $\calC$ to be a collection of functors
  $T_g\colon \calC\to \calC$, $ g\in G$, such that $T_1 = 1$, together
  with functorial isomorphisms
  $$
    \al_{g_1g_2}\colon  T_{g_1}T_{g_2}\simeq T_{g_1g_2}
  $$
  which themselves should satisfy the  compatibility condition: two ways of
  identifying
  $$
    T_{g_1g_2g_3}\simeq T_{g_1}T_{g_2}T_{g_3}
  $$
  should be equal.

  If we have such an action, we define the category $\calC^G$ as the category
  with the following objects and morphisms:
  \begin{itemize}
    \item Objects: an object $X\in \calC$ together with a collection of
    isomorphisms $\ph_g\colon T_G(X)\simeq X$, $g\in G$ such that
    $\ph_1=1_X$ and $\ph_{gh}=\ph_g  T_g (\ph_h)$ (both sides are morphisms
    $T_{gh}(X)\simeq T_{g}T_h(X)\to X$).
    \item Morphisms:
    $$
    \Hom_{\calC^G}(X, X')=\{f\in \Hom_{\calC} (X,X')\st \ph'_g f = \rho_g(f)
    \ph_g\}
    $$
    (FIXME: verify)
  \end{itemize}
  This construction is sometimes called ``equivariantization'', see
  \ocite{etingof-fusion}*{Section 2.7}.
  In particular, if we consider the category $\calC=\vecf$ of finite-dimensional
  vector spaces with trivial action of a group $G$, then $\vecf^G=\Rep G$ is the
  category of finite-dimensional representations of $G$.
\end{example}

In a similar fashion, one can write explicitly a definition of an action of a
group $G$ on a 2-category and definition of $\calC^G$ in this case; one can find
all details in \ocite{hesse}. However, in the case of interest to us, there is
a way around it. Namely, we are only interested in the action of a group on the
$(\infty,0)$ category $\tilde \calC$. By FIXME, every $(\infty, 0)$ category is
a fundamental groupoid of a  topological space: $\calC=\pi_{\le \infty}(X)$.
Thus, it is natural to define the action of $G$ on $\pi_{\le \infty}(X)$ as the
action induced by an action of $G$ on $X$.

Now, we need to define the notion of fixed points. The naive definition of fixed
points of action of $G$ on $X$ is not homotopy invariant, so we cannot use it.
Instead, we will  use the notion of \emph{homotopy fixed point}\index{homotopy fixed point}.

Let $G$ be a topological group (for example, a Lie group). Recall
that then one can define a contractible space $EG$ with a free action of
$G$; the quotient $BG=EG/G$ is called the classifying space of $G$. The name is
justified by the fact that for any space $X$, isomorphism classes of $G$-bundles
on $X$ are in bijection with homotopy equivalence classes of maps $X\to BG$.

\begin{definition}
  Let $X$ be a topological space with a continuous action of a topological group
  $G$. The \emph{homotopy fixed point set} is defined by
  $$
    X^{hG} =\Hom_G(EG, X)
  $$
  where $\Hom_G$ stands for the space of $G$-equivariant continuous maps.
\end{definition}

This definition is rather natural: note that the usual set of $G$-invariant
points can be defined as   $X^G=\Hom_G(pt, X)$, where $pt$ is the one-point set
with trivial action of $G$. One can show that unlike $X^G$, the homotopy fixed
point set $X^{hG}$ is homotopy invariant.

Now we can formulate the oriented version of the cobordism hypothesis.

\begin{theorem}\label{t:cob-hyp-or}
  Let $\calC$ be an $(\infty, n)$ category, and let $\cfd$ be the
  groupoid of fully dualizable objects as in \thref{t:cob-hyp}.
  Then one has a natural action
  of the group $\OO(n)$ on $\cfd$, and we have an equivalence
  $$
    \Fun^{\otimes}(\Bordor, \calC)\simeq (\cfd)^{h\SO(n)}.
  $$
  \end{theorem}
This theorem can be  derived from the (framed) cobordism hypothesis;
details can be found in \ocite{lurie-cobordism}.

\begin{remark}
  One might wonder why we use $\SO(n)$ instead of $\GL_+(n)$. In fact, it does
  not matter: inclusion $\SO(n)\subset \GL_+(n)$ is homotopy equivalence, so the
  set of homotopy fixed points for both actions is the same. However, it is more
  convenient to work with compact groups, hence the use of $\SO(n)$.
\end{remark}




\section{Cobordism hypothesis for $n=2$}\label{s:cobordism-hypothesis-dim2}

Reference: \ocite{hesse}.
$SO(2)$ action on arbitrary 2-category, Serre automorphism, and homotopy fixed
points. Calabi-Yau object.


Main claim:

Fully dualizable objects in $\Alg$ are separable algebras

$\SO(2)$ homotopy fixed points are separable symmetric Frobenius algebras.





Example: $\calC=\Alg$. Symmetric separable Frobenius algebra
